\cotefile{32}
\scanpage{24}

$\underline{\overline{\mathrm{XVIII}}}$. $1.2$

\mspage{6}

\deleted{\unsure{p.} $1.2.3$}

\margin{OK} l. 4 du bas

\unsure{remplacer} $1.2.2.1$ par $(1\unsure{..|.}2.3)$

\smallskip\noindent\rule{.44\linewidth}{.35pt}

\[
\deleted{\mathrm{Lie}\,\mathrm{Pic}_{X}=H^{1}(X,\unsure{o|\mathcal O}_{X})\ ?}
\]
% archive sheet number: 23

\scanpage{25}

\uuline{$\underline{\overline{\mathrm{XVIII}}}$. $1.2$}

\[
\mathrm I\left\{\begin{aligned}
\mathrm I.a\quad & G\longrightarrow p_*p^*G\\
\mathrm I.b\quad & \underline{\mathrm{Hom}}(G_m,G)\xrightarrow{\sim}{\unsure{R|\mathbb R}}^1p_*G
\end{aligned}\right.
\]
\[
\mathrm{II}\ \left\{\ \text{le foncteur $(1.2.1.1)$ est une équivalence.}\right.
\]

\uline{Remarques 1} : \deleted{\unsure{disparaît}} moyennant $(\mathrm I.a)$ et l'hypothèse d'une section
on a une suite exacte
\[
0\longrightarrow H^1(S,G)\longrightarrow {\unsure{\mathbb H|H}}^{1}(\unsure{\mathbb P|P}(\unsure{\mathscr E|\mathcal E|E}),G)
\xrightarrow{\theta}\Gamma(S,R^1p_*G)\longrightarrow0
\]

\uline{Remarque 2} : La pleine fidélité de $(1.2.1.1)$ est locale sur la base et
moyennant la pleine fidélité, la surjectivité essentielle est locale sur la base

\uline{Remarque 3} : $\mathrm I.b$ (et $\mathrm I.a$) entraîne que
$\unsure{\mathrm{Hom}|\underline{\mathrm{Hom}}}(G_m,G)\longrightarrow H^1(\unsure{\underline{P}|\mathbb P|P}(\mathscr E),G)$ est injectif.

\deleted{$\mathrm I)\Rightarrow\mathrm{II})$ pleine fidélité} :

\inserted{\unsure{on se ramène au cas} où il y a une section.}

$\mathrm I\Rightarrow\mathrm{II}$. \uline{pleine fidélité} :

si $K(k_0,\unsure{p|\varphi})\simeq K(k'_0,\varphi')$ alors $\varphi=\varphi'$
\unsure{Alors} $k_0=k'_0$ et $\mathrm I.a$ entraîne la pleine fidélité.

\uline{essentielle surjectivité} : Soit $k\in H^1(\mathbb P(\unsure{E|\mathcal E}),G)$
$\theta(k)=\varphi(\unsure{O|\mathcal O}(1))$ pour un certain $\varphi$ et
$\theta(k)-\varphi(O(1))\in H^1(S,G)$.

$\mathrm{II})\Rightarrow\mathrm I)$ \quad $\mathrm I\text{-}a)$ résulte de la pleine fidélité.

$\mathrm I.b)$ \quad $\deleted{\mathrm{Hom}(G_m,G)\longrightarrow}$ trivial car a \unsure{quasi-inv$^{\mathrm{se}}$|qui précède}.
% archive sheet number: 24

\scanpage{26}

\uline{Lemme} $1.2.\deleted{\unsure{2|3|1|7}.}9$. \uline{Soient} $S$ \uline{un schéma local artinien}, \inserted{$\bar s$ un point géométrique de $S$}, $f:X\longrightarrow S$
\uline{un morphisme propre}\inserted{\uline{et lisse}}\deleted{\unsure{et plat}, de fibres \unsure{géométriques} \deleted{\illegible{1}} \unsure{réduites|connexes}}, \uline{et}
$G$ \uline{un schéma en groupes commutatif lisse de type fini sur} $S$.
\uline{Si} $H^0(X_{\bar s},\unsure{O|\mathcal O})=k(\bar s)$ \uline{et} $\underline{\mathrm{Pic}}^{\tau}X_{\bar s}=0$, \uline{et}
\deleted{alors le foncteur} \uline{si} $G_{\bar s}$ \uline{ne contient pas de sous-groupe} ${\mathbb G}_m$, \uline{alors le} \deleted{le} \uline{foncteur} $f^*$
\uline{de la catégorie des} $G$-\uline{torseurs sur} $S$ \uline{dans celle des} $f^*G$-\uline{torseurs sur} $X$
\uline{est une équivalence.}

% a. below the dividing line; its arrow ends immediately above b.
\inserted{a. Le problème posé est local sur $S$ pour la topologie étale ;
on se ramène donc au cas $S$ strictement hensélien, donc au cas où
\deleted{on peut donc supposer que} $X/S$ admet une section ; le foncteur $f^*$
est alors pleinement fidèle d'après \deleted{$1.2.2.4$}. $1.2.3$ et $(1.2.4.2)$}

b. \uline{Si} $S$ \uline{est le spectre d'un corps}\inserted{\deleted{séparablement clos}} et si $A$ \inserted{est} une variété
abélienne, alors $H^1(X,A)$ est de torsion (Raynaud [1, $\underline{\overline{\mathrm{XIII}}}$ $2.6$]),
donc réunion des images des groupes $H^1(X,{}_nA)$. \deleted{; ceux-ci
sont nuls par $1.2.2.7$}. Si $G$ est extension d'une variété
abélienne $A$ par un groupe \deleted{\unsure{$G_0$}} affine $G_0$, \deleted{$H^1(X,G)$ est donc}
% archive sheet number: 25

\scanpage{27}

\margin{1} $\underline{\overline{\mathrm{XVIII}}}$ $1.2$ \mspage{8}

et si $G^n$ est l'image réciproque de $ {}_nA$, le groupe $H^1(X,G)$
est réunion des groupes $H^1(X,G^n)$. Les groupes $G^n$ sont
affins (non nécessairement lisse) et, sous l'hypothèse faite,
l'assertion résulte de $1.2.\deleted{\unsure{2|1|7}.}8$.

c. Pour $S$ strictement hensélien et $G$ lisse, on a $H^1(S,G)=0$.
il nous faut donc vérifier que $H^1(X,G)=0$. Supposons que
ce soit déjà établi après le changement de base $i:S_0\longrightarrow S$, où
$S_0$ est défini par un idéal de carré nul \inserted{et de longueur $1$} $I$ \deleted{\unsure{du|de}} de $\unsure{O|\mathcal O}_S$. Les
$G$-torseurs sur $X$ sont alors des déformations du $G$-torseur
trivial sur $X_0=X\times_S S_0$. Si $L=e^*\Omega^1_{G/S}$ ($e$ section neutre)
ces déformations sont classifiées par
\[
\deleted{H^1(S_0}
\]
\[
H^1(X_{\deleted{\unsure{0|s}}},f^*\underline{\mathrm{Hom}}(L,I))=0\ ,
\]
et ceci prouve $1.2.\deleted{\unsure{1|2|7}.}9$.
% archive sheet number: 26

\scanpage{28}

$\underline{\overline{\mathrm{XVIII}}}$ $1.2$ \mspage{12}

% Two diagonals cancel the body of the sheet.
\deleted{Pour $S$ artinien, on a $(\mathrm i)\Rightarrow(\mathrm i')$ : si $G$ vérifie $(\mathrm i)$, donc
est fini localement libre, soient $G'$ son dual de Cartier, et
$A(G)$ le schéma en groupe \inserted{(lisse)} des éléments inversibles de l'algèbre
affine de $G'$ ; $G\simeq G''$ s'identifie au sous-groupe de $A(G)$ formé
des caractères de $G'$. Soit donc une suite exacte}
\[
\deleted{0\longrightarrow G\longrightarrow G_1\longrightarrow G_2\longrightarrow0}
\]
\deleted{avec $G_1$ et $G_2$ lisses. La suite exacte longue de cohomologie
fournit, par $1.2.3$. et $1.2.2$ appliqué à $G_1$ et $G_2$}
\[
\deleted{0\longrightarrow R^1p_*G\longrightarrow\mathrm{Hom}(\mathbb G_m,G_1)\longrightarrow\mathrm{Hom}(\mathbb G_m,G_2)}
\]
\deleted{de sorte que $\mathrm{Hom}(\mathbb G_m,G)\xrightarrow{\sim}\deleted{\mathrm{Hom}} R^1p_*G$.}

\deleted{Dans le cas $(\mathrm{iii})$, le problème est local sur $S$ $(1.2.4)$.
Quitte à étendre $k$, on peut supposer l'existence d'une suite
exacte}
\[
\deleted{0\longrightarrow\mathbb G_m^n\longrightarrow G_0\longrightarrow H\longrightarrow0}
\]
\deleted{telle que $H$ ne contienne pas de tore. On applique alors
$1.2.3$ et $1.2.8$ à la suite exacte longue de cohomologie.}
% archive sheet number: 27

\scanpage{29}

% The manuscript is rotated 180 degrees relative to the archive number.
% Three broad diagonal strokes cross the working face.
\deleted{$D$ : $(\deleted{\unsure{V}}$}

\deleted{$\mathbb R^2,\omega$ \quad str. complexe \quad $(V_0,J)$}

\deleted{\unsure{ce|ie|i.e.} $\deleted{V_0}$ \quad $V=V_0\otimes_{\mathbb R}\mathbb C$ et $\omega^{-1}$, \quad $\omega>0$ ?}

\deleted{\uline{\unsure{Grass|Gras}}. \quad \unsure{0|O|$\mathbb C$}.}

\[
\deleted{0\longrightarrow\omega\longrightarrow V\longrightarrow\omega^{-1}\longrightarrow0
\qquad\text{attaché à }D}
\]
\[
\boxed{\deleted{(D,V);V_0,\nabla,0\longrightarrow\omega\longrightarrow V\longrightarrow\omega^{-1}\longrightarrow0,\ \text{str. réelle.}}}
\qquad\deleted{\underline{\mathrm{Aut}}:\mathrm{SL}_2(\mathbb R)}
\]

\deleted{``disque orienté''.}

\deleted{$\mathrm{Aut}(V_0):\mathrm{SL}_2(\mathbb R)/\pm1$}

\deleted{\uline{en part} \quad torseur sous $\mathrm{PSL}_2(\mathbb R)$ (plat) sur disque abst.}

\[
\boxed{\deleted{\longrightarrow\text{torseur plat sous }\mathrm{PSL}_2(\mathbb R)\text{ sur tt courbe}}}
\]
\deleted{(azumaya $\underline{\mathrm{Aut}(V)}$)}

\deleted{a/ se releve-t-il à $\mathrm{SL}_2(\mathbb R)$ ?}

\deleted{$\big[$ fibré en \unsure{dr|dr.} $\pi$ \quad $<$ \quad
\tikz[baseline=(a.base)]{\node[draw,rounded corners=9pt] (a) {parties principales ?};}}

\[
\boxed{\deleted{\deleted{\mathbb P(\unsure{P|p}'(X))}}}
\]
\diagram{\begin{tikzcd}[ampersand replacement=\&,column sep=small,row sep=tiny]
{\deleted{\unsure{g}}}\arrow[d,dash] \& {\deleted{\mathrm{OK}}}\arrow[r,Rightarrow] \& {\deleted{k\mid g}}\arrow[d,dash] \& {\deleted{\mathrm{OK}\ ;\ \Leftrightarrow\text{ pour }k\text{ impair.}}}\\
{\deleted{\unsure{\chi|X}}} \& \& {\deleted{\chi}} \& {}\\
{\tikz[baseline=(a.base)]{\node[draw,circle,inner sep=1pt] (a) {$\deleted{g-1}$};}} \& \& \& {}
\end{tikzcd}}
\deleted{$\mathrm{Aut}(\text{\unsure{rev.|rév.} univ en }x)$}

\diagram{$\begin{array}{lll}
\deleted{\text{obstruction}} & \deleted{H^2(X,\mathbb Z/2)} &
\left[\begin{array}{ll}\deleted{\underline{\text{ou}}}&\deleted{0}\\ \deleted{\underline{\text{ou}}}&\deleted{g-1}\end{array}\right.\\
\deleted{\text{choix}} & \deleted{H^1(X,\mathbb Z/2)} &\\
\deleted{\mathrm{Aut}} & \deleted{H^0(X,\mathbb Z/2)} &
\end{array}$}
\margin{\deleted{$\unsure{\surd|<}$\tikz[baseline=(a.base)]{\node[draw,circle,inner sep=1pt] (a) {$g$};}}}

\deleted{\uline{$\Gamma$ ds $\mathrm{PSL}_2$ peut-il se relever à $\mathrm{SL}_2$} \quad
$x_1\ \unsure{\cdots|\ldots}\ x_{2g}$ avec $\underline{(x_1\,x_2)\ \cdots\ (x_{2g-1}\,x_{2g})=e}$
\quad relever $x_i$ ; \quad \unsure{Rel|rel} $=\pm1$ !}
\[
\deleted{(x_1\,x_2)(x_3\,x_4)=\pm1}
\]
% archive sheet number: 28

\scanpage{30}

\uline{SGA 4 $\underline{\overline{\mathrm{XVII}}}$}

\uline{Vérifier réf} : ``$F$ premier aux car rés. de $X$''
\quad $\underline{\overline{\mathrm{XVII}}}$ $0.13$

Ajouter : $X\longrightarrow S$, \quad $F$ sur $X$ premier aux car de \deleted{\unsure{$X$}} $S$

: ? $\forall s\in S,\quad F\mid X\times_S\mathrm{Spec}(\unsure{O|\mathcal O}_{S,s})$ tel que $p_s$ inv.

\smallskip\hfil\rule{.10\linewidth}{.35pt}\hfil

à \uline{Ajouter} : $1.2.9$ : plat/lisse par \unsure{Illusie|Illusié} --
% archive sheet number: 29

\scanpage{31}

\mspage{L 4}

\uuline{Manque} : \uline{a} \uline{Gysin}\qquad \uline{b} $g^*Rf^!\xrightarrow{\sim}Rf^!g^*$ si \mbox{-\,-\,-}

\uline{c} Un grand nombre de compatibilités \unsure{idiotes|idiots}

\uline{d} Compatibilité sur classes fondamentales locales.

Je n'ai pas l'intention de combler ces manques dans un
temps prévisible. Je me sens toutefois une obligation morale
\margin{$\left[\vphantom{\rule{0pt}{5\baselineskip}}\right.$}de combler $3.1.10.3$., et de donner les compatibilités requises
pour identifier $Rf^!$ à $f^*(\ )[\ ]$ (\unsure{ie|i.e}, telle flèche définie en $3.1$ est telle flèche).

Le risque de faire une vingtaine de pages illisibles et autant de
jours de travail idiot pour moi -- Peut-être le ferais-je avant
que SGA4 ne devienne un livre --

Je ne me sens pas d'obligations morales pour le reste, \unsure{moins|mais}
emmerdant -- Le début de \uline{d} est sans doute à remonter en $1.1$.

\uline{Ci-joint} : tes lettres et notes sur les jacobiennes -- une lettre à Serre --

ta note sur Gysin --

de tes commentaires anciens et nouveaux

une photocopie d'une liste de références dans $\underline{\overline{\mathrm{XVII}}}$.

une liste incomplète de compatibilités manquantes.

une ancienne version, avec tes commentaires, de $3.1$.

Bien à toi\\
P. Deligne\\
% archive sheet number: 30

\scanpage{32}

\uline{Liste de compatibilités manquantes}

a) comp. de ``ch base'' avec $Rf_!\longrightarrow Rf_*$ . \deleted{\unsure{comp}}

b) comp de ``ch base'' avec $f^*(d)[2d]\longrightarrow Rf^!$

c) comp. de \unsure{induction} avec $f^*(d)[2d]\longrightarrow Rf^!$

d) comp. de cl. fond. loc. avec $R{\unsure{g|j}}^!Rf^!=R(fg)^!$

e) classes de correspondance

f) Th de comparaison, sur $\mathbb C$

g) comparaison avec la dualité des courbes par autodualité des jacobiennes

h) morphisme de Gysin et la requête de Jouanolou.

i) ch de base : transp de $(\deleted{\deleted{\unsure{\mathrm{Non}|\mathrm{Var}}}\ Rf_!\,Rg'_*\longrightarrow Rg_*\,Rf'_!}$
\[
R'f^!g^*\mathrel{\blacktriangleleft\mkern-4mu\relbar\mkern-4mu\blacktriangleleft}g'^*Rf^!\quad\text{; cas où }f\text{ est lisse.}
\]
% archive sheet number: 31

\scanpage{33}

\uline{ETAT Avant \uline{Illusie}}

\[
\deleted{H^1(H^0)\longrightarrow H^0(H^1)}
\]
\diagram{\begin{tikzcd}[ampersand replacement=\&,column sep=small,row sep=small]
X \& \underset{(0)}{H^1(H^0)}\arrow[r] \& H^1\arrow[r] \& H^0H^1\arrow[r] \& H^2H^0\\
S \& {} \& H^1\arrow[u]\arrow[ur,two heads] \& {} \& {}
\end{tikzcd}}

% The rest of the writing faces the opposite long edge.
\[
\left\{\begin{aligned}
&\langle\mathcal L,\mathcal M\rangle=\unsure{\mathrm{Tr}|\mathrm N}_{E/S}(\mathcal M)\qquad\text{si }\mathcal L=\unsure{O|\mathcal O}(E)\\
&\text{Tout }F(\mathcal L)\text{ vient d'un }\mathcal M.
\end{aligned}\right.
\]
\[
\longrightarrow O(-F)\longrightarrow O\longrightarrow O_F\longrightarrow0
\]
\[
\begin{aligned}
O(E)\qquad &O(F)\\[-1pt]
&O
\end{aligned}
\qquad
\begin{aligned}
&\langle O(E),\,O(-F))=N_{E/}(O(-F))\,N_{E/}(\ )^{-1}\\
&\det(\ )\det(\ )=\det(\unsure{R|h}f_*\unsure{O_E\otimes O_F|\mathcal O_E\otimes^{L}\mathcal O_F})^{-1}
\end{aligned}
\]
\[
\Big\Vert\langle O(E),O(F)\rangle=\det Rf_*(O_E\otimes O_F)^{-1}
\]
\diagram{\begin{tikzcd}[ampersand replacement=\&,row sep=small,column sep=small]
\langle 1\ 1\rangle\arrow[rr,dash,"{\scriptstyle :}"']\arrow[d,dash] \& {} \& \langle1\ 1\rangle\arrow[d,dash]\\
\langle \unsure{O|0|\mathcal O}X\ OY\rangle\arrow[d,dash] \& \underline{\langle f,f^{-1}\rangle} \& \langle OY\ OX\rangle\arrow[d,dash]\\
O \& {} \& O
\end{tikzcd}}
% archive sheet number: 32

\scanpage{34}

\[
\begin{aligned}
X'_1\quad&Y_1\qquad X_1\quad T\\
X_2\quad&Y_2
\end{aligned}
\qquad\qquad
\begin{aligned}
&\underline{\text{\unsure{desc. et.|desc. ét.|desc. el.}}}\\
&\text{desc. }\underline{\text{finie}}.
\end{aligned}
\]
\diagram{\begin{tikzcd}[ampersand replacement=\&,row sep=small,column sep=small]
{} \& \unsure{x|\times} \& x \& {\unsure{\mathrm{Sym}|\mathrm{Symm}|\mathrm{Supp}}\,T'}\arrow[d]\\
X_1\cup Y_2 \& (-1)^{v_x f\,v_x g} \& (g^{v_x f}/f^{v_x g}) \& {\mathrm{Sym}\,T}\\
{} \& {\mathrm{Sym}\,T''} \& {\mathrm{Sym}\,T} \& {}
\end{tikzcd}}
\[
\begin{aligned}
X_1\cup T&\qquad(-1)^{v_x f\,v_x g}\quad g^{v_x}\qquad{}^{:}\\
X_{\unsure{1|4}}\cup Y_2&\qquad A_x\longrightarrow A\qquad\equiv h\unsure{,|.}
\end{aligned}
\]
\diagram{\begin{tikzcd}[ampersand replacement=\&,column sep=small,row sep=small]
{} \& T^{\unsure{\prime\prime|\wedge|1}}\arrow[r,shift left=1.2]\arrow[r,shift right=1.2] \& T'\arrow[r] \& {\deleted{\unsure{S}}T}\arrow[dd]\\
{\deleted{X_1\cup T}} \& {X_1\quad Y_1} \& {} \& {}\\
{\deleted{T\cup Y_2}} \& {X_1\quad T} \& {X_3\quad Y_3} \& S\\
{} \& {X_2\quad Y_2} \& {} \& {}
\end{tikzcd}}
\diagram{\begin{tikzcd}[ampersand replacement=\&,column sep=small,row sep=small]
M_2 \& M_1 \& M_0 \& {}\\
A_{\unsure{2|1}}\arrow[r,shift left=3.6]\arrow[r]\arrow[r,shift right=3.6] \&
A_1\arrow[l,shift left=1.8]\arrow[l,shift right=1.8]\arrow[r,shift left=2.4]\arrow[r] \&
A_0\arrow[l,shift left=2.4]\arrow[r] \& A
\end{tikzcd}}
\[
X_1\ Y_3\mid X_3\ Y_2\qquad
\begin{aligned}
&\underline{\text{ops}}:\quad\underline{A_0\longleftarrow A}\\
&\overset{?}{\Rightarrow}\quad\text{ops}\qquad A_0=A
\end{aligned}
\]
% archive sheet number: 33

\scanpage{35}

ETAT

\uline{PRIMITIF.}

\uline{contient} : ? compatibilité

\unsure{commentaire|commutation} de G.

% Very faint traces below the title, possibly transferred from the reverse.
% crop-candidate: scan 35
\figure{\begin{tikzpicture}[x=.62cm,y=.62cm,draw=black!32,line width=.25pt]
\draw (.14,4.68) -- (.32,4.50);
\draw[->] (1.95,.79) -- (3.23,.79);
\draw (8.9,1.34) -- (9.12,1.37);
\draw (1.9,7.4) -- (3.7,7.4);
\draw[<-] (1.9,7.3) -- (3.5,7.3);
\draw[->] (7.2,6.45) -- (8.45,6.45);
\draw[->] (5.85,3.55) -- (5.9,5.15);
\draw[->] (3.75,3.1) -- (5.9,3.1);
\draw[<->] (1.8,2.25) -- (3.4,2.3);
\draw[<-] (7.35,2.3) -- (8.5,2.3);
\draw (3.9,8.8) -- (9.3,2.2);
\draw (1.4,5.8) -- (5.6,.3);
\draw (8.9,5.7) -- (9.6,5.7) -- (9.3,4.9);
\foreach \x/\y in {1.4/8.2,2.2/7.1,2.45/6.7,3.4/7.0,4.1/5.5,4.5/5.43,6.2/5.5,7.75/4.35,2.25/2.65,3.95/2.93,4.4/3.05,6.2/3.05,9.0/1.16}
  {\draw (\x,\y) ellipse (.075 and .10);}
\end{tikzpicture}}
% archive sheet number: 34

\scanpage{36}

% Long, oblique strokes cancel the working sheet.
\diagram{\begin{tikzcd}[ampersand replacement=\&,column sep=small,row sep=small]
{\deleted{\mathrm{Cat}}} \& {}\arrow[dd,thick,"{\deleted{f^{\bullet}_!}}"'] \& {\boxed{\deleted{{}^{*}\mathcal A\qquad{}^{*}\unsure{\mathcal B|B}}}} \& {}\\
{} \& {} \&
{\tikz[baseline=(a.base)]{\node[draw,circle,inner sep=1pt] (a) {\deleted{OK}};}} \& {}\\
{} \& {} \& {\boxed{\deleted{C_{*}\qquad \unsure{\mathcal D|D}_{*}}}} \& {}\arrow[uu,thick,"{\deleted{?}}"']
\end{tikzcd}}
\diagram{\begin{tikzcd}[ampersand replacement=\&,column sep=large]
{\deleted{*}}\arrow[r,thick]\arrow[d,"{\deleted{!}}" near start,"{\deleted{?}}" near end,"{\deleted{!}}"' near end] \& {\deleted{\ast}}\arrow[d]\\
{\deleted{\underset{0}{\times}}}\arrow[r,thick] \& {\deleted{\underset{1}{\times}}}
\end{tikzcd}}
\diagram{\begin{tikzcd}[ampersand replacement=\&,column sep=large,row sep=large]
{\deleted{\times}}\arrow[r,"{\deleted{\mathrm{Hom}({\unsure{g|f}}^{*}F,G)}}"] \& {\deleted{\times}}\\
{\deleted{\times}}\arrow[u,"{\blacktriangle}"{description,pos=.06,inner sep=0pt}]\arrow[r] \&
{\deleted{\times}}\arrow[u,"{\blacktriangle}"{description,pos=.06,inner sep=0pt},"{\deleted{\mathrm{Hom}({}_{!}\ ,{}^{!}\ )}}"']
\end{tikzcd}}

\deleted{ou}

\diagram{\begin{tikzcd}[ampersand replacement=\&,column sep=large,row sep=large]
{\deleted{{}^{\unsure{B|F}}\times}}\arrow[d,thick,"{\deleted{F}}"' near start,"{\deleted{\mathrm{Hom}(f_!F,G)}}"',"{\deleted{G}}"' near end]
\& {\deleted{\times}}\arrow[l]\arrow[d,"{\deleted{\mathrm{Hom}(F,f^!G)}}"]\\
{\deleted{\times}} \& {\deleted{\times}}\arrow[l,"{\deleted{G}}" near start,"{\deleted{F}}" near end,"{\deleted{\mathrm{Hom}(f^*G,F)}}"]
\end{tikzcd}}
\diagram{\fbox{\begin{tikzcd}[ampersand replacement=\&,column sep=small,row sep=small]
{\deleted{\bullet}}\arrow[d,"{\deleted{f_!}}"'] \& {\deleted{\odot}}\arrow[l,"{\deleted{g^*}}"']\arrow[d,"{\deleted{f_!}}"] \&
{\deleted{\bullet}}\arrow[r,"{\deleted{g_*}}"] \& {\deleted{\bullet}} \& {}\\
{\deleted{\odot}} \& {\deleted{\bullet}}\arrow[l,"{\deleted{g^*}}"]\arrow[ul,phantom,"{\tikz[baseline=-.4ex]{\draw[->,thick] (.14,0) .. controls (.14,.08) and (.07,.14) .. (0,.14);}}" description] \&
{\deleted{\bullet}}\arrow[u,"{\deleted{f^!}}"]\arrow[r,"{\deleted{g_*}}"'] \& {\deleted{\bullet}}\arrow[u,"{\deleted{f^!}}"']\arrow[ul,phantom,"{\tikz[baseline=-.4ex]{\draw[->,thick] (.14,0) .. controls (.14,.08) and (.07,.14) .. (0,.14);}}" description]
\& {\deleted{\underline{\text{\unsure{dér|dév|dém}}}:\underline{\text{fib}+\text{cofib}}}}
\end{tikzcd}}}
\deleted{$\longrightarrow$ \uline{ch base} : \uline{2 \unsure{descr|desc.|dessin|dessins}}.}

\diagram{\begin{tikzcd}[ampersand replacement=\&,column sep=small,row sep=small]
{\deleted{\underline{\text{adj}}:}} \& {\deleted{\mathrm{Hom}(f_!K,L)}}\arrow[d,thick] \& {\deleted{\sim}} \& {\deleted{\mathrm{Hom}(K,f^!L)}}\arrow[d,thick] \& {\deleted{/X}}\\
{} \& {\deleted{\mathrm{Hom}_V(f_!^V K,L)}} \& {\deleted{\sim}} \& {\deleted{\mathrm{Hom}(K,f_V^!L)}} \& {\deleted{/V}}
\end{tikzcd}\quad\deleted{$\underline{\simeq}$\uline{\unsure{préf|pref|prof}}}}
\[
\deleted{Rf_*\,R\mathrm{Hom}(f_!K,L)\ \sim}
\]
% archive sheet number: 35

